\documentclass{situationist-preprint}

\title{Hobbes on the Problem, Locke on the Solution}
\msid{POL-05-a}
\msfamily{POL-05}
\msdiscipline{Political Philosophy}
\mstype{Principal manuscript}
\mskeywords{state of nature; Hobbes; Locke; assurance problem; sovereignty; limited government; right of resistance; legitimacy}

\begin{document}
\maketitle

\begin{abstract}
Hobbes and Locke both justify government by comparison with a state of nature, and they disagree about what that state is like and what follows. Hobbes's state of nature is a war of every one against every one, and its remedy is an absolute sovereign; Locke's is a state of liberty governed by natural law, whose inconveniences are remedied by a limited government held in trust. We argue that Hobbes correctly identifies the problem and that Locke provides the better account of legitimate government. The problem is the absence of a common authoritative judge and enforcer, and it does not depend on Hobbes's psychology: it arises among agents of limited altruism who are uncertain of each other's intentions and judge in their own cases. Locke's three inconveniences are the same problem, and his own text concedes that they lead to war. But Hobbes's solution reproduces the problem between sovereign and subject, as Locke's objection about the polecats and the lion shows, while Locke's limited, fiduciary government solves it without reproducing it. Locke's account of how government is legitimated, by consent, fails; his account of what legitimate government is survives, restated as a condition of justifiability.
\end{abstract}

\section{The question}

The device of the state of nature asks what human life would be like without political authority, and uses the answer to say what political authority is for and what it may do. \citet[ch.~13]{Hobbes1651} answered that it would be a condition of war, in which there is no place for industry, agriculture, navigation or letters, and the life of man is solitary, poor, nasty, brutish and short. \citet[\S\S 4--15]{Locke1689} answered that it would be a state of perfect freedom and equality, bounded by a law of nature that obliges everyone and that everyone may execute, and that its defect is not war but a set of inconveniences, arising from each person's being judge in her own case. From these different starting points they drew different conclusions. Hobbes concluded that the remedy is a sovereign whose power is unlimited and undivided, against whom subjects have no right of resistance except in self-preservation \citep[chs.~18, 21]{Hobbes1651}. Locke concluded that the remedy is a government limited by the purposes for which it is instituted, holding its power in trust, and liable to be resisted and dissolved when it betrays that trust \citep[\S\S 149, 221--222]{Locke1689}.

The question who was right is usually treated as a single question. We argue that it is two. The first is what problem political authority solves: what is wrong with the state of nature. The second is what account of legitimate government follows: what form of authority solves the problem rightly. Our thesis is that Hobbes gives the better answer to the first and Locke to the second. The combination is not an eclectic compromise. Locke's account of legitimate government is better precisely because it solves the problem Hobbes identified, which Hobbes's own solution does not.

We do not adjudicate the interpretive literature; we rely on the rational-choice reconstructions of Hobbes by \citet{Gauthier1969}, \citet{Kavka1986} and \citet{Hampton1986}, and on \citet{Simmons1992}'s reading of Locke. Nor do we claim that either state of nature describes a historical condition; both thinkers used the device to identify the function of government, and our claims concern the function.

\section{Framework}

\begin{definition}\label{d:son}
A \emph{state of nature} is a condition of persons among whom there is no common power with authority to make rules binding on all of them, to judge disputes under the rules, and to enforce its judgements.
\end{definition}

\begin{definition}\label{d:problem}
The \emph{problem of order} is the problem of securing, among persons in a state of nature, the assurance that others will refrain from attacking them and will respect their holdings, sufficient for each to pursue her plans without devoting her resources to defence.
\end{definition}

\begin{definition}\label{d:account}
An \emph{account of legitimate government} specifies the form and limits of authority under which a government's rule is legitimate, and the conditions under which that legitimacy is lost.
\end{definition}

We assume three facts about persons, each weaker than its Hobbesian counterpart.

\begin{assumption}[Vulnerability]\label{a:vuln}
Persons are roughly equal in their vulnerability: each can be gravely harmed by others, singly or in combination.
\end{assumption}

\begin{assumption}[Limited altruism]\label{a:altru}
Most persons give some weight to the interests of others, but more to their own and their dependants'; some give little or none; and no one can reliably tell, of a stranger, which kind of person he is.
\end{assumption}

\begin{assumption}[Partiality of judgement]\label{a:partial}
Persons disagree in good faith about what they are owed, and each tends to judge disputes in which she is a party in her own favour.
\end{assumption}

Assumption~\ref{a:vuln} is Hobbes's natural equality \citep[ch.~13]{Hobbes1651}, without the claim that the weakest can kill the strongest by stealth in every case. Assumption~\ref{a:altru} replaces psychological egoism by what \citet[ch.~2]{Kavka1986} calls predominant egoism, combined with uncertainty about types. Assumption~\ref{a:partial} is the core of Locke's own diagnosis: it is unreasonable for men to be judges in their own cases, since self-love will make them partial to themselves and their friends \citep[\S 13]{Locke1689}.

\section{Hobbes on the problem}

\begin{proposition}[Assurance]\label{p:assure}
Under Assumptions~\ref{a:vuln} and~\ref{a:altru}, persons in a state of nature cannot secure the assurance of Definition~\ref{d:problem}, even if most of them would prefer mutual restraint to mutual predation.
\end{proposition}

\begin{proof}
Consider two persons who each prefer mutual restraint to any other outcome, but who each prefer defensive preparation, or a pre-emptive strike, to being the only one who restrains himself. Their situation is a stag hunt: mutual restraint is an equilibrium, but so is mutual defence, and each should choose restraint only if confident the other will. By Assumption~\ref{a:altru}, each knows that some persons are predators, and cannot tell of the other whether he is one. By Assumption~\ref{a:vuln}, the cost of misplaced trust is grave. A person who restrains himself against someone who does not risks death; a person who prepares against someone who restrains himself loses only the cost of the preparation. Under such uncertainty, preparation is the prudent choice even for those who prefer restraint, and since each knows this of the others, the expectation of preparation is self-confirming. This is the reasoning \citet[ch.~13]{Hobbes1651} attributes to diffidence, and it requires only that some be predators and that no one can tell who.
\end{proof}

Proposition~\ref{p:assure} is the core of Hobbes's argument as \citet{Kavka1986} and \citet{Hampton1986} reconstruct it, and \citet{Skyrms2004} has shown how robust the stag hunt's risk-dominant equilibrium is. The war of all against all does not require bad persons, only that good persons be unable to recognize each other with confidence when the cost of error is grave.

\begin{lemma}[Locke's inconveniences]\label{l:locke}
Under Assumption~\ref{a:partial}, the three inconveniences Locke attributes to the state of nature produce the problem of order.
\end{lemma}

\begin{proof}
Locke names three wants: of an established, settled, known law; of a known and indifferent judge; and of a power to back and support a right sentence \citep[\S\S 124--126]{Locke1689}. Suppose a dispute arises over a holding or an injury. Without a settled law, the parties may disagree in good faith about what the law of nature requires; by Assumption~\ref{a:partial}, each will judge the dispute in her own favour. Without an indifferent judge, there is no one whose verdict both must accept. Each, being the executor of the law of nature, may then enforce her own verdict, and the other, judging the enforcement unjust, may resist and retaliate. The result is the condition Locke calls the state of war, which he says is one great reason of men's putting themselves into society: to avoid it, where there is no authority to appeal to, the least difference is apt to end in it \citep[\S 21]{Locke1689}. Since no one can be assured that her holdings and person will be respected when every dispute may escalate in this way, the problem of order arises.
\end{proof}

\begin{theorem}\label{t:problem}
The problem that political authority solves is the absence of a common authoritative judge and enforcer, and it arises among persons of the kind described by Assumptions~\ref{a:vuln} to~\ref{a:partial}, whatever their moral character. Hobbes's identification of the problem is correct, and Locke's is a special case of it.
\end{theorem}

\begin{proof}
By Proposition~\ref{p:assure}, the problem of order arises from uncertainty and vulnerability alone; by Lemma~\ref{l:locke}, it arises from partiality alone. In both cases what is missing is the feature named in Definition~\ref{d:son}: a common power with authority to make, judge and enforce rules. Hobbes identifies this lack as the source of war \citep[ch.~13]{Hobbes1651}. Locke identifies it as the source of inconveniences and concedes that they lead to war. His state of nature differs from Hobbes's not in the problem but in its severity, which he understates: he supposes that a law of nature known to all would moderate disputes, and Assumption~\ref{a:partial} shows that a law known to all, but interpreted by each in her own favour, does not.
\end{proof}

Theorem~\ref{t:problem} does not vindicate Hobbes's psychology, nor his claim that in the state of nature nothing is unjust. Locke is right that there are natural rights and duties before government, and POL-18 argues that rights are real normative constraints, not creatures of law. What Hobbes saw is that the existence of such rights does not solve the problem, since the problem lies in the absence of a common power to determine and enforce them, not in the absence of the rights themselves.

\section{Locke on the solution}

Hobbes inferred from the problem that its solution must be a sovereign with unlimited and undivided power. His argument was that any limitation of the sovereign requires a judge of whether the limits have been exceeded; if that judge is the sovereign, there is no limit, and if it is another body or the subjects, then there is a second power, and divided power is a return to the state of nature \citep[chs.~18, 29]{Hobbes1651}.

\begin{lemma}[Reproduction]\label{l:repro}
An unlimited sovereign reproduces the problem of order between itself and its subjects.
\end{lemma}

\begin{proof}
The sovereign is a person or body of persons, and Assumptions~\ref{a:vuln} to~\ref{a:partial} apply to it as to anyone. Between sovereign and subject there is, on Hobbes's account, no common judge: the sovereign judges its own case. By Assumption~\ref{a:altru}, subjects cannot be assured that the sovereign is not a predator; by Assumption~\ref{a:vuln}, its power to harm them is far greater than any private person's; and by Assumption~\ref{a:partial}, it will judge disputes between itself and them in its own favour. The subjects therefore lack, with respect to the sovereign, the assurance of Definition~\ref{d:problem}. This is Locke's objection: to think that men would guard against the mischiefs done by polecats or foxes, and be content, indeed think it safety, to be devoured by lions \citep[\S 93]{Locke1689}. An absolute monarch, Locke concludes, is in the state of nature with respect to those under his dominion \citep[\S\S 90--91]{Locke1689}.
\end{proof}

\begin{lemma}[Limitation without division]\label{l:limit}
The problem of order can be solved by a government whose power is limited by standing law, separated among bodies that check one another, and held in trust for the purposes for which it was instituted.
\end{lemma}

\begin{proof}
Hobbes's argument assumes that a limitation requires a single further judge who would then be a rival sovereign. It does not. A government whose legislative, executive and judicial functions are assigned to distinct bodies, each bound by standing law and each able to check the others, has a common judge for every dispute between private persons, since the courts decide them under settled law, and it has a procedure, not a rival person, for disputes over the limits of power. Locke's three wants are met: there is a settled law, an indifferent judge, and a power to enforce \citep[\S\S 124--131]{Locke1689}. And because no single body is unchecked, Lemma~\ref{l:repro} does not apply: subjects have, with respect to each part of the government, a judge other than that part. Hobbes's own evidence, the civil war of the 1640s, shows that a division of power without settled rules for disputes between its parts is unstable; constitutional governments that supplied such rules have persisted for centuries without the war he predicted of divided sovereignty.
\end{proof}

\begin{theorem}\label{t:solution}
Locke's account of legitimate government, on which government is limited, held in trust, and liable to be resisted when it betrays its trust, is better than Hobbes's account, on which the sovereign is unlimited and unresisted.
\end{theorem}

\begin{proof}
A government is a solution to the problem of order only if it secures the assurance of Definition~\ref{d:problem} for those it governs. By Lemma~\ref{l:repro}, Hobbes's sovereign fails to secure it with respect to the sovereign itself. By Lemma~\ref{l:limit}, Locke's government secures it without reproducing the problem. Locke's residual right of resistance, which Hobbes held to be a return to war, is the condition under which the assurance holds: a government that knows it can be resisted when it betrays its trust has a reason not to betray it that an unlimited sovereign lacks. An account of legitimate government that solves the problem is better than one that does not, and Locke's is therefore better.
\end{proof}

The comparison can also be made through the criterion of POL-02, on which a state is legitimate if and only if it can be justified to every reasonable person subject to it. A reasonable person who accepts government in order to escape the insecurity of the state of nature cannot be given sufficient reason, from her own standpoint, to accept a sovereign who stands to her as a lion, since that would leave her as insecure as before, against a more powerful threat. She can be given sufficient reason to accept a limited, accountable government. Hobbes's sovereign therefore fails the criterion and Locke's government can meet it.

This restatement also marks where Locke goes wrong. His account of how government is legitimated is consent, express or tacit \citep[\S\S 95--122]{Locke1689}, and \citet{Hume1748} showed that tacit consent by residence is not consent that could bind. Locke's account of \emph{what} legitimate government is, limited, fiduciary and accountable, survives this failure, since the features that make it the right form of government are those that make it justifiable to those it governs.

\section{Objections}

\begin{objection}[Who shall be judge?]
Locke's right of resistance requires that someone judge whether the government has betrayed its trust. Locke answers that the people shall be judge \citep[\S 240]{Locke1689}. But the people are not a single judge; they are many persons who will judge in their own favour, by Assumption~\ref{a:partial}. The right of resistance therefore reinstates the state of nature whenever a faction believes itself wronged.
\end{objection}

\begin{reply}
Locke anticipated the objection and answered it with a threshold: the people are not moved to resist by small mismanagement or particular wrongs, but only by a long train of abuses that makes a design against them visible to all \citep[\S 225]{Locke1689}. Resistance is the collective judgement of a people, not the private judgement of a faction, and the conditions for it are rarely met. POL-23 argues that revolution is justified only under conditions of extreme injustice, and that its justification depends on outcomes that cannot be known in advance; that uncertainty is a further reason for the threshold. A residual right of collective resistance under extreme conditions does not reinstate the state of nature in ordinary conditions. The alternative, a sovereign that cannot be resisted whatever it does, guarantees that the subjects remain in the state of nature with respect to it.
\end{reply}

\begin{objection}[Hobbes's psychology is false]
Hobbes derived the war of all against all from a psychology on which every voluntary act aims at the agent's own good. The psychology is false; most people are not egoists. So Hobbes did not identify the problem correctly; he identified a problem that would arise among beings unlike us.
\end{objection}

\begin{reply}
Proposition~\ref{p:assure} does not use psychological egoism. It uses Assumption~\ref{a:altru}, on which most people give some weight to others but no one can tell of a stranger how much. That is the predominant egoism of \citet{Kavka1986}, and it is a weaker and more plausible premise than Hobbes's. \citet{Hampton1986} argues that Hobbes's argument is driven less by egoism than by the conflict generated by disagreement and the desire for glory, which our Assumption~\ref{a:partial} captures. The problem survives the rejection of Hobbes's psychology, and that is why we credit Hobbes with identifying it rather than with explaining its source.
\end{reply}

\begin{objection}[Stateless societies]
Anthropologists have described stateless societies in which order is maintained without a common power, by kinship, reciprocity and reputation. Hobbes's war of all against all is not their condition, and Locke's milder picture is closer to the truth.
\end{objection}

\begin{reply}
The mechanisms such societies rely on, repeated interaction among people who know each other and can observe each other's conduct, are mechanisms that reduce the uncertainty of Assumption~\ref{a:altru}. They work where the population is small and stable and fail where it is large and mobile, which is why they do not refute Proposition~\ref{p:assure} for societies of the kind that political philosophy is about. POL-07 argues, moreover, that stateless societies are coherent but unstable: their protective associations tend to consolidate into state-like institutions, which is the Hobbesian problem working itself out.
\end{reply}

\section{Conclusion}

Hobbes saw that what is wrong with the state of nature is not that people are wicked, but that without a common power to judge and enforce, even the well-disposed cannot be assured of one another, and Locke's inconveniences turn out, on his own account, to be that problem. Hobbes's solution, however, places over the subjects a power with respect to which they remain in the state of nature. Locke's limited, accountable government solves the problem Hobbes identified without reproducing it, and it is the better account of legitimate government, even if its own account of how government is legitimated must be replaced.

\begin{thebibliography}{99}\small
\bibitem[Gauthier(1969)]{Gauthier1969} Gauthier, D.~P. (1969). \emph{The Logic of Leviathan: The Moral and Political Theory of Thomas Hobbes}. Clarendon Press.
\bibitem[Hampton(1986)]{Hampton1986} Hampton, J. (1986). \emph{Hobbes and the Social Contract Tradition}. Cambridge University Press.
\bibitem[Hobbes(1651)]{Hobbes1651} Hobbes, T. (1651). \emph{Leviathan}. Cited by chapter.
\bibitem[Hume(1748)]{Hume1748} Hume, D. (1748). Of the original contract. In \emph{Essays, Moral, Political, and Literary}.
\bibitem[Kavka(1986)]{Kavka1986} Kavka, G.~S. (1986). \emph{Hobbesian Moral and Political Theory}. Princeton University Press.
\bibitem[Locke(1689)]{Locke1689} Locke, J. (1689). \emph{Two Treatises of Government}. Cited by section of the Second Treatise.
\bibitem[Simmons(1992)]{Simmons1992} Simmons, A.~J. (1992). \emph{The Lockean Theory of Rights}. Princeton University Press.
\bibitem[Skyrms(2004)]{Skyrms2004} Skyrms, B. (2004). \emph{The Stag Hunt and the Evolution of Social Structure}. Cambridge University Press.
\end{thebibliography}

\end{document}
